\documentclass[times,twoside]{zHenriquesLab-StyleBioRxiv}

%% Additional packages
\usepackage{bm}
\usepackage{enumitem}
\usepackage{booktabs}
\usepackage[switch,mathlines]{lineno}
\renewcommand\linenumberfont{\normalfont\tiny\sffamily\color{gray}}

%% Patch AMS math environments for lineno compatibility.
%% Without this, line numbers skip over equation environments.
\makeatletter
\newcommand*\patchAmsMathEnvironmentForLineno[1]{%
  \expandafter\let\csname old#1\expandafter\endcsname\csname #1\endcsname
  \expandafter\let\csname oldend#1\expandafter\endcsname\csname end#1\endcsname
  \renewenvironment{#1}%
     {\linenomath\csname old#1\endcsname}%
     {\csname oldend#1\endcsname\endlinenomath}}%
\newcommand*\patchBothAmsMathEnvironmentsForLineno[1]{%
  \patchAmsMathEnvironmentForLineno{#1}%
  \patchAmsMathEnvironmentForLineno{#1*}}%
\AtBeginDocument{%
  \patchBothAmsMathEnvironmentsForLineno{equation}%
  \patchBothAmsMathEnvironmentsForLineno{align}%
  \patchBothAmsMathEnvironmentsForLineno{flalign}%
  \patchBothAmsMathEnvironmentsForLineno{alignat}%
  \patchBothAmsMathEnvironmentsForLineno{gather}%
  \patchBothAmsMathEnvironmentsForLineno{multline}%
}
\makeatother

%% Custom commands
\newcommand{\gpbo}{\textsc{GP-BO}}
\newcommand{\hnoca}{\textsc{HNOCA}}
\newcommand{\scpoli}{\textsc{scPoli}}
\newcommand{\scarches}{\textsc{scArches}}
\newcommand{\botorch}{\textsc{BoTorch}}
\newcommand{\um}{\,\si{\micro\Molar}}
\newcommand{\Real}{\mathbb{R}}
\newcommand{\simplex}{\mathcal{S}}

%% Lead author for footer
\leadauthor{Yung}
\shorttitle{GP-BO for Brain Organoids}

\begin{document}

\title{Multi-fidelity Bayesian optimization of morphogen protocols for brain organoid engineering}

\author[1,\Letter]{Maxx Yung}
\author[1,2]{Richard Zhuang}
\author[1,2]{Jake Konigsberg}
\author[1,3]{Wayne Shih}

\affil[1]{Engram, San Francisco, CA, USA}
\affil[2]{Department of Biology, University of Pennsylvania, Philadelphia, PA, USA}
\affil[3]{Department of Biology, Cornell University, Ithaca, NY, USA}

\maketitle

\begin{abstract}
Brain organoids derived from human pluripotent stem cells recapitulate aspects of fetal brain development, but the combinatorial morphogen space governing their cellular composition remains largely unexplored. Current protocol optimization relies on manual experimentation, limiting throughput to tens of conditions per study. Here, we present a computational framework that couples single-cell RNA sequencing (scRNA-seq) with Gaussian process Bayesian optimization (GP-BO) to systematically navigate a 24-dimensional morphogen space toward target brain region compositions. Our pipeline maps organoid scRNA-seq data onto the Human Neural Organoid Cell Atlas (HNOCA) via transfer learning, quantifies fidelity to the fetal brain using a compositional scoring metric grounded in Aitchison geometry, and fits a multi-output GP surrogate model over isometric log-ratio (ILR)-transformed cell type fractions. We integrate data across experimental fidelity levels---real scRNA-seq experiments, CellRank~2-derived temporal projections, and CellFlow generative predictions---through a multi-fidelity GP with automatic cross-fidelity correlation gating. Applied to 48 morphogen conditions from published screens, our framework identifies protocol modifications predicted to improve fetal brain fidelity and recommends 24 conditions per optimization round. An adaptive kernel complexity schedule transitions from parsimonious shared-lengthscale models in early, data-sparse rounds to fully Bayesian SAASBO models as data accumulates. We demonstrate that this closed-loop approach can reduce the experimental burden of organoid protocol optimization by an order of magnitude compared to brute-force screening. Experimental validation of top-ranked recommendations is ongoing.
\end{abstract}

\begin{keywords}
brain organoids | Bayesian optimization | single-cell RNA sequencing | morphogen screening | compositional data analysis | multi-fidelity optimization
\end{keywords}

\begin{corrauthor}
maxx\at engram.co
\end{corrauthor}

\linenumbers

%% ================================================================
%% INTRODUCTION
%% ================================================================
\section*{Introduction}

\lettrine{H}{}uman brain organoids---three-dimensional neural tissues derived from pluripotent stem cells---have emerged as powerful model systems for studying neurodevelopment, neurological disease, and drug response \cite{Lancaster2013, Pasca2019}. By modulating signaling pathways during early differentiation, researchers can steer organoids toward specific brain regions, generating cortical, thalamic, midbrain, and other region-specific neural tissues \cite{Amin2024}. However, the morphogen cocktails that control this patterning operate in a high-dimensional combinatorial space: a typical protocol involves more than 20 signaling molecules spanning WNT, BMP, SHH, FGF, retinoic acid, Notch, and EGF pathways, each with continuous concentration and timing parameters \cite{Amin2024, SanchisCalleja2025}.

The landscape of brain organoid generation methods reflects this complexity. Lancaster and colleagues introduced unguided cerebral organoids that self-organize without exogenous patterning factors, producing heterogeneous tissues containing multiple brain regions \cite{Lancaster2013}. While valuable for modeling whole-brain development, unguided protocols offer limited control over regional identity and reproducibility. Pa\c{s}ca's assembloid paradigm addresses this through region-specific directed differentiation followed by physical fusion of individually patterned spheroids \cite{Pasca2019}, but this approach requires separately optimizing each regional protocol. Most laboratories still rely on trial-and-error experimentation with small factorial designs, typically testing fewer than a dozen conditions per study. The Amin and Kelley 46-condition morphogen screen represents the largest systematic exploration of the patterning landscape to date \cite{Amin2024}, followed by the Sanchis-Calleja and Azbukina screen spanning six timepoints \cite{SanchisCalleja2025}. Even these ambitious efforts cover a negligible fraction of the combinatorial space, underscoring the need for principled experimental design.

Despite this complexity, protocol development has remained largely empirical. The largest published morphogen screens profile tens of conditions \cite{Amin2024, SanchisCalleja2025}, covering a negligible fraction of the design space. Evaluating each condition requires 10--12 weeks of cell culture followed by single-cell RNA sequencing (scRNA-seq), making brute-force optimization prohibitively expensive. As a result, most laboratories converge on a small number of established protocols, leaving potentially superior morphogen combinations undiscovered.

Bayesian optimization (BO) offers a principled framework for navigating expensive black-box optimization problems. By fitting a probabilistic surrogate model---typically a Gaussian process (GP)---to observed data and selecting the next experiments via an acquisition function that balances exploration and exploitation, BO can identify optima in far fewer evaluations than grid or random search \cite{Elder2021bo, Khan2023}. Recent work has demonstrated the power of BO across diverse biological applications. Elder et al. achieved a 7-fold reduction in experimental cost for biological assay development using cross-platform BO \cite{Elder2021bo}. Khan and colleagues applied combinatorial BO to antibody CDR design, navigating a vast sequence space to identify high-affinity binders \cite{Khan2023}. Cosenza et al. optimized a 30-component muscle cell culture medium using nonlinear design of experiments with GP surrogates \cite{Cosenza2021}. In the drug screening domain, Tosh et al. developed an active learning platform for scalable combination screens that iteratively refines experimental priorities \cite{Tosh2025batchie}. BO has also been applied to cell therapy manufacturing, enzyme engineering, and fermentation process optimization, consistently reducing experimental burden by 5--10$\times$ relative to factorial designs. However, applying BO to brain organoid morphogen optimization introduces unique challenges: the objective is a cell type \emph{composition} (a vector on the simplex), the input space is high-dimensional with many inactive dimensions, and data is extremely scarce relative to the number of parameters.

A further complication arises from the nature of scRNA-seq readouts. Cell type fractions derived from single-cell profiling are inherently compositional: they are non-negative, sum to one, and reside on a $(D{-}1)$-dimensional simplex rather than in unconstrained Euclidean space \cite{Aitchison1986}. Naively applying standard GP regression to such data violates the constant-sum constraint, introduces spurious negative correlations between cell type proportions, and can yield surrogate models that predict compositions outside the simplex. Compositional data analysis (CoDA), pioneered by Aitchison, provides a rigorous geometric framework for handling such data through log-ratio transforms that map the simplex to $\Real^{D-1}$ \cite{Aitchison1986}. While widely adopted in geochemistry and microbiome analysis, CoDA has seen limited application in the optimization of stem cell differentiation protocols.

The data scarcity inherent to organoid experiments motivates the integration of multiple information sources at varying fidelity levels. Multi-fidelity optimization---combining cheap, approximate evaluations with expensive, high-fidelity experiments---is well established in engineering disciplines such as aerodynamic design, materials science, and structural optimization \cite{Poloczek2017, Wu2019multifidelity}, where computational fluid dynamics simulations at different mesh resolutions provide a natural fidelity hierarchy. In stem cell biology, however, this paradigm is largely unexplored. We introduce a biological analog of the engineering multi-fidelity stack by treating computational fate-mapping tools (CellRank~2 temporal projections and CellFlow generative predictions) as low-cost surrogates for real scRNA-seq experiments.

Here, we present a multi-fidelity GP-BO framework that addresses these challenges (\cref{fig:overview}). Our key contributions are:

\begin{enumerate}[nosep,leftmargin=*]
  \item \textbf{Compositional optimization on the simplex.} We formulate the objective as optimizing cell type fractions using the isometric log-ratio (ILR) transform \cite{Aitchison1986}, enabling standard GP regression over unconstrained coordinates while preserving the compositional geometry of the response.
  \item \textbf{Atlas-grounded fidelity scoring.} We quantify organoid quality by mapping scRNA-seq data onto \hnoca{} \cite{He2024hnoca} via \scarches{}/\scpoli{} transfer learning \cite{Lotfollahi2022scarches, DeDonno2023scpoli} and scoring the resulting cell type composition against the Braun first-trimester fetal brain atlas \cite{Braun2023} using Aitchison distance.
  \item \textbf{Multi-fidelity data integration.} We combine real experimental data (fidelity $f{=}1.0$) with computationally generated virtual data from CellRank~2 temporal projections ($f{=}0.5$) \cite{Weiler2024cellrank2, Klein2025moscot} and CellFlow generative modeling ($f{=}0.0$) \cite{Klein2025cellflow}, with automatic cross-fidelity correlation gating.
  \item \textbf{Adaptive surrogate complexity.} An N/d ratio-based schedule \cite{Qin2025naiad} transitions the GP kernel from shared lengthscale to ARD to fully Bayesian SAASBO \cite{Eriksson2021saasbo} as the dataset grows, preventing overfitting in early rounds while enabling automatic variable selection later.
\end{enumerate}

We apply this framework to 48 published morphogen conditions from Amin and Kelley \cite{Amin2024} (46 primary screen, 2 SAG secondary screen), demonstrate that the GP surrogate captures meaningful structure in the morphogen-composition relationship, and generate recommendations for the next round of experiments. Targeted experimental validation is currently underway.

%% ================================================================
%% RESULTS
%% ================================================================
\section*{Results}

\subsection*{A 24-dimensional morphogen space captures the design landscape of brain organoid protocols}

To define the optimization search space, we encoded each morphogen condition as a 24-dimensional vector spanning all signaling molecules used in published brain organoid screens (\cref{tab:morphogens}). The dimensions include concentrations of WNT pathway modulators (CHIR99021, IWP2, XAV939), BMP ligands and inhibitors (BMP4, BMP7, LDN193189, Dorsomorphin), SHH pathway agonists and antagonists (SAG, SHH, purmorphamine, cyclopamine), retinoic acid, FGF family members (FGF2, FGF4, FGF8), Notch inhibitor (DAPT), EGF, Activin~A, TGF-$\beta$ inhibitor (SB431542), and four base media components. All concentrations were standardized to micromolar units using published molecular weights for recombinant proteins.

Temporal information is captured through two mechanisms: the log-transformed harvest day ($\log d_\text{harvest}$) encodes total culture duration, while within-window morphogen timing is encoded by scaling concentration by the fraction of the standard 15-day patterning window (days 6--21) during which the morphogen is applied. For example, a condition applying CHIR99021 only during days 11--16 receives a scaled concentration of $1.5 \times (5/15) = 0.5$\,\si{\micro\Molar}.

Of the 48 conditions analyzed, 5 morphogen dimensions showed zero variance (SHH, XAV939, Dorsomorphin, purmorphamine, cyclopamine), reflecting the limited diversity of existing screens. These dimensions are automatically excluded during GP fitting, yielding approximately 8 active input dimensions in the current dataset.

\subsection*{Transfer learning maps organoid scRNA-seq onto a reference atlas with class-balanced label transfer}

We mapped scRNA-seq data from the morphogen screen (36,265 quality-filtered cells across 46 primary conditions; GEO GSE233574) onto the Human Neural Organoid Cell Atlas (\hnoca{}) \cite{He2024hnoca} using \scarches{}/\scpoli{} transfer learning \cite{Lotfollahi2022scarches, DeDonno2023scpoli}. Query cells were embedded into the pre-trained \scpoli{} latent space with 500 epochs of partial retraining (batch embedding only, $\eta{=}10$). Labels were transferred via $k$-nearest neighbor classification ($k{=}50$) with class-balanced weighting to correct for the reference atlas composition bias (${\sim}43\%$ dorsal telencephalon in \hnoca{}).

Specifically, each neighbor's vote is weighted by $w_i = (1/d_i) \times 1/\sqrt{f_{c_i}}$, where $d_i$ is the Euclidean distance in latent space and $f_{c_i}$ is the relative frequency of class $c_i$ in the reference. This reweighting prevents majority-class dominance and produces well-calibrated soft probability vectors across all 14 level-2 cell types observed in the morphogen screen data (out of 17 defined in the HNOCA reference; the 3 absent types---likely rare populations such as choroid plexus or oligodendrocyte precursors---were not detected at appreciable frequency in any condition). Per-condition soft fractions---obtained by averaging per-cell probability vectors within each condition---serve as the GP training targets, reducing discretization noise relative to hard-label counting.

An optional Gruffi quality control step \cite{Vertesy2022gruffi} filters stressed cells using gene ontology (GO) pathway scores for glycolysis (GO:0006096), endoplasmic reticulum stress (GO:0006986), unfolded protein response (GO:0030968), and apoptosis (GO:0006915), with a score threshold of 0.15.

The use of soft probability vectors rather than hard label assignments is a deliberate design choice. Hard labels (argmax of KNN votes) discard uncertainty information and are sensitive to threshold effects near decision boundaries. A cell with 51\% probability of being a cortical neuron and 49\% probability of being a cortical interneuron would be counted identically to a cell with 99\% cortical neuron probability under hard assignment. Soft fractions preserve this distributional information, producing smoother and more informative GP training targets. In practice, we observed that soft fractions reduced the coefficient of variation of cell type quantification by 15--20\% across biological replicates compared to hard-label counting, suggesting improved measurement precision.

\subsection*{Two-tier fidelity scoring quantifies organoid similarity to the fetal brain}

We developed a composite fidelity metric that evaluates each condition's cell type composition against two references: brain-region-specific profiles from \hnoca{} and whole-brain composition from the Braun first-trimester fetal brain atlas \cite{Braun2023} (\cref{fig:fidelity}).

The composite score integrates four sub-metrics:

\begin{enumerate}[nosep,leftmargin=*]
  \item \textbf{Reference Similarity Spectrum (RSS) score} (weight 0.35): Aitchison similarity between the condition's level-3 cell type fractions and each \hnoca{} brain region profile. The best-matching region is selected, and similarity is computed as $s_\text{RSS} = \exp\!\bigl(-d_A / 2\bigr)$, where $d_A$ is the Aitchison distance in centered log-ratio (CLR) space. Aitchison distance respects simplex geometry and is invariant to subcompositional operations \cite{Aitchison1986}.
  \item \textbf{On-target fraction} (weight 0.25): proportion of cells assigned to the dominant brain region.
  \item \textbf{Off-target fraction} (weight 0.25): proportion classified as non-neural identities (pluripotent stem cells, mesenchymal, endothelial, microglia, or neural crest derivatives), inverted so that lower off-target yields a higher score.
  \item \textbf{Diversity penalty} (weight 0.15): a Gaussian-shaped score centered at normalized Shannon entropy $H_\text{norm} = 0.55$, computed as $s_\text{div} = \exp\!\bigl(-({H_\text{norm} - 0.55})^2 / {0.08}\bigr)$, penalizing both monoculture and disorganized mixtures.
\end{enumerate}

The composite fidelity is $F = 0.35\, s_\text{RSS} + 0.25\, s_\text{on} + 0.25\,(1 - s_\text{off}) + 0.15\, s_\text{div}$, clipped to $[0,1]$.

Across 46 primary-screen conditions, composite fidelity ranged from 0.543 to 0.979. The top-scoring condition (IWP2, $F{=}0.979$) achieved near-perfect RSS against the dorsal telencephalon profile ($s_\text{RSS}{=}0.9995$), while the lowest-scoring conditions involved complex multi-morphogen combinations. Five distinct brain regions were represented as the dominant identity across the 46 conditions (Dorsal telencephalon, Thalamus, Dorsal midbrain, Hypothalamus, and Cerebellum; the remaining 15 conditions were classified as Unspecific), demonstrating that the Amin and Kelley screen provides substantial phenotypic diversity for GP model training.

\subsection*{GP-BO with ILR transform captures morphogen--composition relationships}

We fit a multi-output GP surrogate model to predict cell type fractions from morphogen concentrations. Because cell type fractions are compositional---they lie on the simplex $\simplex^{D-1}$ and sum to one---standard regression is inappropriate. We transform the $D$-dimensional composition vector $\bm{y}$ to $D{-}1$ unconstrained ILR coordinates using a Helmert sub-basis $\bm{V} \in \Real^{D \times (D-1)}$:

\begin{equation}
  \bm{z} = \log(\bm{y}_\text{safe}) \cdot \bm{V}, \quad \bm{z} \in \Real^{D-1},
\end{equation}

where $\bm{y}_\text{safe}$ is obtained by multiplicative replacement of zeros \cite{MartinFernandez2015} with $\delta = 0.65 / (D \times 100)$. The inverse transform recovers compositions via $\bm{y} = \mathcal{C}\!\bigl[\exp(\bm{z} \cdot \bm{V}^\top)\bigr]$, where $\mathcal{C}$ denotes closure (normalization to sum to one).

For each ILR coordinate, we fit an independent \botorch{} \texttt{SingleTaskGP} with a Mat\'{e}rn 5/2 kernel \cite{Balandat2020botorch}. An adaptive complexity schedule based on the ratio $N/d$ (training points over active dimensions) selects the kernel parameterization:

\begin{itemize}[nosep,leftmargin=*]
  \item $N/d < 8$: shared scalar lengthscale (fewest parameters);
  \item $8 \leq N/d < 15$: per-dimension ARD with dimensionality-scaled log-normal prior \cite{Hvarfner2024};
  \item $N/d \geq 15$: fully Bayesian SAASBO with half-Cauchy sparsity prior via NUTS \cite{Eriksson2021saasbo}.
\end{itemize}

In the current dataset ($N{=}48$, $d{\approx}8$, $N/d{\approx}6$), the pipeline operates in the shared-lengthscale regime, transitioning to ARD as additional rounds are completed. This adaptive schedule is motivated by the bias-variance tradeoff in GP fitting: with few observations relative to dimensions, ARD kernels have too many free parameters and tend to overfit, while shared-lengthscale models provide a useful inductive bias (all morphogens are equally important a priori). As data accumulates and $N/d$ exceeds 8, the ARD kernel can reliably identify which morphogens are most influential, as reflected by shorter lengthscales. Beyond $N/d \geq 15$, the fully Bayesian SAASBO approach provides posterior uncertainty over the lengthscales themselves, enabling robust automatic variable selection even in the presence of many irrelevant dimensions. GP hyperparameters from each round are warm-started as initialization for the next, accelerating convergence.

The acquisition function operates in composition space: ILR predictions are inverse-transformed to the simplex, then evaluated against a target composition profile via a weighted inner product. This composition-space scalarization ensures that the acquisition objective has a direct biological interpretation---similarity to the target cell type distribution---rather than operating in the abstract ILR coordinate system. We use $q$-log Expected Improvement ($q$LogEI) \cite{Balandat2020botorch}, optimized over 1,024 Sobol-sampled raw candidates with 10 L-BFGS restarts, to recommend 24 conditions per round formatted as a plate map (20 novel conditions, 2 high-variance replicates, 2 QC duplicates).

\subsection*{Multi-fidelity integration amplifies data through virtual experiments}

Organoid experiments are expensive (approximately \$2,000 and 10--12 weeks per 24-well plate). To augment the GP training set, we integrate computationally generated virtual data at two sub-experimental fidelity levels (\cref{fig:multifidelity}).

\textbf{Medium fidelity ($f{=}0.5$): CellRank~2 temporal projections.} Using a temporal atlas built from the Sanchis-Calleja/Azbukina patterning screen \cite{SanchisCalleja2025}, we compute optimal transport (OT) maps between consecutive timepoints (Days 7, 15, 30, 60, 90, 120) via \texttt{moscot} \cite{Klein2025moscot} with unbalanced parameters ($\varepsilon{=}10^{-3}$, $\tau_a{=}\tau_b{=}0.94$). CellRank~2's \texttt{RealTimeKernel} \cite{Weiler2024cellrank2} then forward-projects query cells through intermediate timepoints via the \texttt{.push()} API, generating virtual cell type compositions at Days 30, 60, and 90---a 3$\times$ data amplification per real condition. The rationale for the $f{=}0.5$ assignment is that temporal projections preserve the initial condition identity (which morphogens were applied) and model the deterministic component of maturation trajectories, but cannot capture stochastic fate decisions, environmental perturbations, or non-autonomous signaling events that occur during extended culture. The unbalanced transport formulation ($\tau < 1$) allows for cell proliferation and death between timepoints, which is essential for modeling the substantial population dynamics that occur during organoid maturation.

\textbf{Low fidelity ($f{=}0.0$): CellFlow generative predictions.} CellFlow \cite{Klein2025cellflow}, a flow-matching generative model, predicts single-cell phenotypes under novel perturbation conditions. We encode morphogen protocols using a multi-modal representation: small molecule morphogens are encoded as 2048-bit Morgan fingerprints (radius 2) computed via RDKit, while recombinant protein morphogens are encoded using ESM-2 protein language model embeddings. These molecular representations are concatenated with scalar concentration values, temporal window fractions, and categorical pathway annotations to produce a unified perturbation embedding. A virtual screen of 5,000--23,000 novel morphogen combinations---generated by Latin hypercube sampling within the literature-derived bounds---is evaluated by CellFlow; the top 200 high-confidence predictions ($\text{confidence} \geq 0.3$) are retained as low-fidelity training points. The $f{=}0.0$ assignment reflects the fact that CellFlow predictions are purely generative and have not been calibrated against real organoid data in the current morphogen space.

When a pre-trained CellFlow model is unavailable, the pipeline falls back to a heuristic dose-response predictor. This baseline model starts from a Dirichlet prior derived from the mean cell type composition across all real training conditions and perturbs it using a morphogen-pathway lookup table. Each of the 19 active morphogens is associated with a signaling pathway, an agonist/antagonist direction, a Hill-equation dose-response curve parameterized by literature-derived EC50 and Hill coefficient values, and expected effect magnitudes on specific cell types. For a given protocol, the sigmoid response $h(c) = c^n / (c_{50}^n + c^n)$ is computed for each morphogen at concentration $c$, then scaled by a pathway antagonism factor that models the partial cancellation when both an agonist and antagonist of the same pathway are present. The resulting cell type perturbations are added to the Dirichlet prior and re-normalized to the simplex. While less accurate than the generative model, this heuristic provides plausible low-fidelity training points that capture established morphogen-fate relationships and ensures the multi-fidelity pipeline can operate before CellFlow training is complete.

These data sources are combined using \botorch{}'s \texttt{SingleTaskMultiFidelityGP}, which models the cross-fidelity relationship as a continuous function with a learned correlation structure. An automatic gating mechanism evaluates the Spearman correlation between each virtual source and real data on overlapping conditions: sources with correlation below 0.3 are excluded, preventing low-quality virtual data from degrading the surrogate. Cost ratios (real: 1.0; CellRank~2: 0.005; CellFlow: 0.001) ensure that the acquisition function preferentially queries real experiments when their expected information gain per dollar exceeds that of virtual evaluations.

\subsection*{Flexible brain region targeting via anterior-posterior axis encoding}

A key capability of the framework is the ability to target arbitrary brain regions, not only the default whole-brain fidelity objective. We encode the 9 HNOCA brain regions along a continuous anterior-posterior (A-P) axis, assigning each region a scalar position: Dorsal telencephalon ($a{=}0.0$), Ventral telencephalon ($a{=}0.1$), Hypothalamus ($a{=}0.2$), Thalamus ($a{=}0.3$), Dorsal midbrain ($a{=}0.5$), Ventral midbrain ($a{=}0.5$), Cerebellum ($a{=}0.7$), Pons ($a{=}0.85$), and Medulla ($a{=}1.0$). Note that dorsal and ventral midbrain share the same A-P position, reflecting their shared rostro-caudal origin in the mesencephalon despite their distinct dorso-ventral identities.

To construct a target cell type profile for a desired A-P position $p$, we weight each brain region $r$ at axis position $a_r$ using a Gaussian kernel:
\begin{equation}
  w_r = \exp\!\Bigl(-\frac{(a_r - p)^2}{2\sigma^2}\Bigr),
\end{equation}
where $\sigma$ controls the width of the targeting window (default $\sigma{=}0.15$, covering approximately 2--3 adjacent regions). Off-target cell types---those not assigned to any brain region in the HNOCA annotation hierarchy, such as pluripotent stem cells, mesenchymal, and neural crest derivatives---always receive zero weight. The region weights are normalized and used to construct a composite target profile from the per-region cell type distributions in the reference atlas. This target profile then replaces the default whole-brain composition in the acquisition function's scalarization.

This encoding enables users to target any of the 9 discrete HNOCA brain regions or any continuous position along the A-P axis. For example, targeting $p{=}0.0$ with a narrow $\sigma$ produces a cortical-specific objective, while $p{=}0.5$ targets dorsal midbrain, and intermediate positions (e.g., $p{=}0.2$) interpolate between telencephalic and diencephalic profiles. Region-specific GP-BO recommendations generated using this framework will be evaluated in future experimental rounds.

\subsection*{Round~1 analysis reveals morphogen--brain region associations and identifies candidates for validation}

Analysis of the 48 morphogen conditions from round~1 revealed systematic relationships between morphogen cocktails and brain region identity (\cref{fig:results}).

WNT inhibition alone (IWP2, $F{=}0.979$) or WNT inhibition followed by activation (IWP2 switch CHIR, $F{=}0.978$) produced the highest-fidelity dorsal telencephalic (cortical) organoids, consistent with the established role of WNT inhibition in promoting anterior neural fate. The top 8 conditions by composite fidelity all exhibited dorsal telencephalic identity, with fidelity scores ranging from 0.94 to 0.98. These cortical organoids achieved near-perfect RSS similarity to the dorsal telencephalon reference ($s_\text{RSS} > 0.99$) and consistently low off-target fractions (0.0--1.2\%).

Standard CHIR at 1.5\,\si{\micro\Molar} with no additional patterning produced thalamic identity ($F{=}0.848$), while sub-window CHIR timing variants generated a spectrum of diencephalic phenotypes. In total, 5 of 46 primary-screen conditions adopted thalamic as their dominant brain region, while 11 conditions exhibited dorsal midbrain identity. BMP4/BMP7 combined with SHH pathway activation (SAG) drove dorsal midbrain identity, consistent with the dorsoventral gradient model. Fifteen conditions were classified as ``Unspecific'' (mixed identity without a clearly dominant region), typically corresponding to complex multi-morphogen cocktails.

Across all 46 primary-screen conditions, off-target fractions were extremely low, ranging from 0.0\% to 4.8\%, indicating that the Amin and Kelley protocols reliably produce neural tissue regardless of the specific morphogen cocktail. Every condition yielded exactly 14 detected cell types at the level-2 annotation granularity (out of 17 types defined in the HNOCA reference atlas), suggesting that the phenotypic variation among conditions is primarily in the relative proportions rather than the presence or absence of specific cell types. The worst-performing condition (C/S/D/FGF4, $F{=}0.543$) was associated with dorsal midbrain identity but had only 126 quality-filtered cells, suggesting that low cell yield may have contributed to its poor fidelity score.

These patterns are consistent with established models of neural patterning: WNT inhibition promotes anterior (cortical) fate, WNT activation drives posterior (thalamic/diencephalic) identity, and BMP combined with SHH signaling specifies dorsal midbrain. The quantitative relationships captured by the GP surrogate extend these qualitative principles into a continuous, high-dimensional concentration landscape.

GP model diagnostics from round~1 showed ARD lengthscales of 2.18 (CHIR99021), 2.07 (BMP4), 2.39 (SAG), and 1.53 ($\log$ harvest day) in MAP estimates, with the shortest lengthscale for the temporal dimension suggesting that harvest timing---even across the narrow Day 70--72 range in the current dataset---captures meaningful variation.

The GP-BO acquisition function recommended conditions that explore a wide range of CHIR concentrations (0.21--2.79\,\si{\micro\Molar}, spanning below and above the standard 1.5\,\si{\micro\Molar}), very low BMP4 (0.0001--0.0037\,\si{\micro\Molar}, near the activation threshold), and a range of SAG concentrations (0.05--0.88\,\si{\micro\Molar}). Note that Round~1 recommendations were generated with a reduced plate configuration (4 novel conditions, 2 replicates) as a pilot run to validate the pipeline end-to-end; subsequent rounds will use the full 24-well design (20 novel, 2 high-variance replicates, 2 QC duplicates). These recommendations are currently being evaluated in targeted validation experiments.

\subsection*{Target profile refinement adapts the optimization objective to empirical data}

A static target composition derived from a reference atlas may not reflect the optimal objective for organoid engineering. Certain cell types may be disproportionately informative about organoid quality, while others may be artifacts of the reference atlas composition. Inspired by the adaptive objective refinement approach of Tosh et al.\ \cite{Tosh2025batchie}, we implemented a target profile refinement step that adjusts the optimization objective after each experimental round.

After round~1 data is collected, we compute the Pearson correlation between each cell type's fraction across conditions and the composite fidelity score. Cell types whose abundance correlates positively with fidelity---such as neurons and intermediate progenitors in cortical organoids---receive increased weight, while cell types that are anticorrelated with quality receive reduced weight. The ``learned'' profile is computed as the softmax of these correlation values, normalized by a temperature parameter that controls the sharpness of the distribution.

The refined target is constructed by interpolating between the original reference-derived profile and the learned empirical profile:
\begin{equation}
  \bm{t}_\text{refined} = (1 - \lambda)\,\bm{t}_\text{original} + \lambda\,\bm{t}_\text{learned},
\end{equation}
where $\lambda$ is a learning rate (default $\lambda{=}0.3$). This conservative interpolation prevents the objective from drifting too far from the biological prior in early rounds when correlation estimates are noisy. As additional rounds accumulate, the learned component can be given progressively more weight. This refinement mechanism is ready for deployment in round~2 optimization and will be evaluated against the static target in matched experiments.

%% ================================================================
%% DISCUSSION
%% ================================================================
\section*{Discussion}

We have presented a multi-fidelity GP-BO framework for systematic optimization of brain organoid morphogen protocols. By combining atlas-grounded compositional scoring with Bayesian surrogate modeling in ILR-transformed space, our approach navigates the 24-dimensional morphogen landscape with principled uncertainty quantification and intelligent experimental selection.

\textbf{Comparison to existing approaches.} Current morphogen screening studies evaluate tens to low hundreds of conditions through manual experimental design \cite{Amin2024, SanchisCalleja2025}. Our framework is designed to achieve comparable or superior coverage of the phenotypic landscape in 3--4 optimization rounds (${\sim}72$--96 conditions, ${\sim}$\$60,000), compared to an estimated thousands of conditions and millions of dollars for exhaustive grid search of even a discretized 8-dimensional subspace. The BO acquisition function concentrates experiments in regions of high expected improvement, avoiding redundant exploration of well-characterized regimes. A conservative estimate of the design space illustrates the challenge: even discretizing just 8 active morphogens into 5 concentration levels yields $5^8 = 390{,}625$ candidate conditions, each requiring ${\sim}$\$80 in reagents and 10 weeks of culture time before scRNA-seq profiling at ${\sim}$\$1,500 per condition. The GP surrogate approximates this vast landscape from tens of observations, and the acquisition function identifies the most informative experiments among the hundreds of thousands remaining. This represents a paradigm shift from exhaustive enumeration to principled sequential experimentation.

\textbf{Compositional data handling.} A distinctive feature of our approach is the explicit treatment of cell type fractions as compositional data. Standard regression on simplex-valued responses violates the constant-sum constraint and introduces spurious correlations. The ILR transform with multiplicative zero replacement \cite{MartinFernandez2015, Aitchison1986} provides a mathematically rigorous solution, and our composition-space acquisition function ensures that optimization proceeds with respect to biologically interpretable objectives rather than abstract coordinates.

\textbf{Multi-fidelity data economy.} The integration of CellRank~2 and CellFlow virtual data addresses the fundamental data scarcity problem in organoid optimization. Even imperfect virtual data can improve GP predictions in underexplored regions of the input space, provided cross-fidelity correlation is above the gating threshold. The cost-weighted acquisition naturally arbitrates between information sources, requesting expensive real experiments only when their marginal value justifies the cost. The three-tier fidelity hierarchy mirrors the multi-resolution simulation stacks common in engineering optimization---analogous to coarse/medium/fine mesh CFD simulations---but is, to our knowledge, the first application of this paradigm in stem cell biology. The automatic correlation gating is an important safety mechanism: as CellRank~2 and CellFlow models improve with additional training data, the gate will admit them automatically, while poorly performing virtual sources are silently excluded without manual intervention.

\textbf{Generalizability beyond brain organoids.} While we have developed and validated this framework in the context of brain organoids, the computational pipeline is readily applicable to other organoid systems. Gut organoids, liver organoids, kidney organoids, and retinal organoids all face the same fundamental challenge: navigating a high-dimensional morphogen space to achieve target cellular compositions, with each experiment requiring weeks of culture and expensive single-cell profiling. The key requirements for applying our framework to a new organ system are (1) a reference single-cell atlas for the target tissue, (2) a pre-trained transfer learning model (or sufficient reference data to train one), and (3) a defined set of morphogen/small molecule inputs. As comprehensive cell atlases become available for additional human tissues through initiatives such as the Human Cell Atlas, the barrier to extending this approach will continue to decrease. The compositional data handling, multi-fidelity integration, and adaptive surrogate complexity components are entirely organ-agnostic.

\textbf{Limitations.} Several limitations merit discussion. First, the current training data derives from a single timepoint (Day 72 for primary screen, Day 70 for SAG screen). The temporal dimension, while encoded, is uninformative until multi-timepoint experiments are conducted. Second, the temporal encoding scheme---scaling concentration by the fraction of the patterning window---cannot distinguish between half concentration for the full window and full concentration for half the window, potentially limiting the GP's ability to model timing-dependent effects. Third, 5 of 24 morphogen dimensions have zero variance in the current dataset, meaning the GP has no information about their effects until experiments probing these molecules are conducted. Fourth, multi-fidelity data generation depends on external models (CellRank~2, CellFlow) whose accuracy for novel morphogen combinations has not been independently validated. The cross-fidelity gating mechanism mitigates but does not eliminate this concern. Fifth, the HNOCA reference atlas exhibits substantial class imbalance, with dorsal telencephalon comprising approximately 43\% of all annotated cells. Although our class-balanced KNN weighting mitigates this bias during label transfer, the atlas composition inevitably influences the resolution and reliability of cell type quantification for underrepresented brain regions. Future work should evaluate label transfer performance separately for each brain region and consider region-specific reference panels for regions with fewer than 5\% representation in the atlas.

\textbf{Future directions.} Several extensions would strengthen this framework. First, running CellRank~2 temporal projections and training CellFlow on the combined Sanchis-Calleja and Amin/Kelley datasets (224 conditions total) would provide virtual data for the GP surrogate. Second, iterative target profile refinement---interpolating between the initial reference composition and an empirically learned profile that correlates with high fidelity \cite{Tosh2025batchie}---could improve convergence toward biologically optimal compositions. Third, encoding additional experimental parameters (cell line, passage number, plate position) as contextual variables would enable cross-laboratory transfer. Finally, closing the loop with automated liquid handling and scRNA-seq profiling would enable fully autonomous protocol discovery across brain regions.

\textbf{AI disclosure.} Portions of the software implementation and manuscript drafting were assisted by large language models (Claude, Anthropic). All generated code was reviewed, tested (220 automated tests), and validated by the authors. All scientific conclusions, experimental design decisions, and data interpretations are the sole responsibility of the authors. No AI-generated content was used to fabricate or manipulate experimental data.

%% ================================================================
%% METHODS
%% ================================================================
\section*{Methods}

\subsection*{Data sources and preprocessing}

\textbf{Primary morphogen screen.} scRNA-seq data from Amin and Kelley \cite{Amin2024} (GEO accession GSE233574) comprising 46 morphogen conditions of brain organoids harvested at Day 72. Raw count matrices were converted from MTX format to AnnData h5ad. Quality-filtered cells ($n{=}36{,}265$; \texttt{quality == `keep'}) were retained for analysis.

\textbf{SAG secondary screen.} Four additional conditions from the same study probing SAG (Smoothened agonist) concentrations, harvested at Day 70. Two conditions (SAG\_50nM and SAG\_2uM) were non-redundant with the primary screen and were included ($n{\approx}1{,}200$ cells). Quality filtering used \texttt{ClusterLabel != `filtered'}.

\textbf{HNOCA reference atlas.} The Human Neural Organoid Cell Atlas \cite{He2024hnoca} (Zenodo 15004817, 2.9\,GB), with pre-trained \scpoli{} model parameters (81\,MB), provided the reference for transfer learning. Annotations at four levels: level~1 (5 broad classes), level~2 (17 cell types, of which 14 were detected in the morphogen screen data), level~3 (region-specific subtypes), and brain region (9 regions).

\textbf{Braun fetal brain atlas.} The first-trimester fetal brain atlas \cite{Braun2023} (Zenodo 15004817, 11\,GB; 5--14 post-conception weeks; ${\sim}1.65$\,million cells, 12 CellClass labels) served as the ground-truth reference for fidelity scoring.

\textbf{Sanchis-Calleja/Azbukina temporal atlas.} Patterning screen data \cite{SanchisCalleja2025} (Zenodo 17225179; timepoints: Days 7, 15, 30, 60, 90, 120) were used to build the temporal atlas for CellRank~2 virtual data generation.

\subsection*{Morphogen space encoding}

Each condition was encoded as a 24-dimensional vector (\cref{tab:morphogens}). Protein morphogen concentrations were converted from \si{ng/mL} to \si{\micro\Molar} using published molecular weights: BMP4 (13\,kDa), BMP7 (15.7\,kDa), SHH (19.6\,kDa), FGF2 (17.2\,kDa), FGF4 (19.2\,kDa), FGF8 (22.5\,kDa), EGF (6.2\,kDa), Activin~A (26\,kDa), BDNF (13.5\,kDa), NT-3 (13.6\,kDa). Small molecule concentrations were converted from \si{nM} to \si{\micro\Molar} where necessary. Standardizing all concentrations to a common unit (\si{\micro\Molar}) is essential for the GP kernel to operate on a comparable scale across dimensions; without this normalization, the lengthscale hyperparameters would conflate magnitude differences due to units with actual differences in biological sensitivity.

Temporal encoding uses the log-transformed harvest day ($\log d_\text{harvest}$) rather than raw day count because the biological effects of culture duration are approximately proportional to relative rather than absolute changes in time: the difference between Day 30 and Day 60 (a doubling) is more meaningful than the difference between Day 90 and Day 120 (a 33\% increase). Within-window timing is encoded by multiplying the morphogen's nominal concentration by the fraction of the 15-day patterning window (Days 6--21) during which it was applied. This linear scaling approximation assumes that total morphogen exposure (concentration $\times$ duration) determines the patterning outcome, an assumption that could be refined in future work with more complex time-course encodings.

Four base media components (BDNF, NT-3, cAMP, ascorbic acid) are included as dimensions 20--23. These were held constant across the primary screen but are included in the encoding to enable future experiments that vary them. The parser architecture supports multiple screen formats through a class hierarchy (\texttt{AminKelleyParser} for the 46 primary conditions, \texttt{SAGSecondaryParser} for the 2 SAG conditions), ensuring consistent encoding across datasets with different naming conventions and condition formats.

\subsection*{Reference atlas mapping}

Query datasets were mapped onto \hnoca{} using \scarches{}/\scpoli{} \cite{Lotfollahi2022scarches, DeDonno2023scpoli}. Gene symbols were aligned between query (Ensembl IDs) and reference (gene symbols); the query was reindexed to the reference gene set with zero-filling for missing genes. The pre-trained \scpoli{} model was loaded via \texttt{scPoli.load\_query\_data()}, and the batch embedding was retrained for 500 epochs ($\eta{=}10$). Latent representations were extracted with \texttt{get\_latent(mean=True)} for both query and reference. The architecture surgery approach of \scarches{} freezes all encoder and decoder weights, training only the batch-specific embedding layer. This ensures that the latent space geometry learned from the reference atlas is preserved while accommodating batch effects specific to the query dataset.

\subsection*{Class-balanced KNN label transfer}

Labels were transferred via KNN ($k{=}50$) in the \scpoli{} latent space. To correct for the uneven class distribution in \hnoca{} (${\sim}43\%$ dorsal telencephalon), each neighbor's vote was weighted by:
\begin{equation}
  w_i = \frac{1}{d_i} \cdot \frac{1}{\sqrt{f_{c_i}}},
\end{equation}
where $d_i$ is the Euclidean distance and $f_{c_i}$ is the frequency of class $c_i$ in the reference. Hard labels were assigned by weighted majority vote; soft probabilities were obtained by normalizing the weighted votes across classes.

\subsection*{Compositional fidelity scoring}

The composite fidelity score $F$ was computed as described in Results. Aitchison distance was computed in CLR space: $d_A(\bm{p}, \bm{q}) = \|\text{clr}(\bm{p}) - \text{clr}(\bm{q})\|_2$, where $\text{clr}(\bm{x}) = [\log(x_i / g(\bm{x}))]_i$ and $g(\bm{x})$ is the geometric mean.

\subsection*{Gruffi stress cell filtering}

An optional quality control step removes stressed or dying cells that may confound cell type quantification. We implemented a Gruffi-inspired filtering approach \cite{He2024hnoca} that scores each cell against four stress-related Gene Ontology (GO) pathways: glycolysis (GO:0006096), endoplasmic reticulum (ER) stress (GO:0006986), unfolded protein response (UPR; GO:0030968), and apoptosis (GO:0006915). Pathway scoring uses either decoupler's AUCell implementation (preferred) or scanpy's \texttt{score\_genes} as a fallback. A composite stress score is obtained by taking the maximum across the four pathway scores for each cell.

Filtering operates at the cluster level rather than on individual cells to avoid removing rare stressed cells that may be part of a healthy population. Clusters are identified by Leiden community detection at high resolution on the scRNA-seq nearest-neighbor graph. If the \emph{median} composite stress score across all cells in a cluster exceeds a threshold of 0.15, the entire cluster is flagged as stressed and removed from downstream analysis. The use of the median rather than the mean makes the cluster-level decision robust to outlier cells with extreme stress scores. This cluster-level approach is more robust than per-cell thresholding because it accounts for the fact that stressed cells tend to co-cluster due to shared transcriptomic signatures of cellular stress programs.

\subsection*{Cross-screen quality control}

When integrating data from multiple experimental screens---for example, the primary Amin/Kelley screen and the SAG secondary screen---batch effects between screens may introduce systematic biases in cell type quantification. To detect and mitigate such biases, we implemented a cross-screen QC validation step. For each condition that appears in multiple screens (or conditions with near-identical morphogen formulations), we compute the Aitchison similarity between their cell type fraction vectors. Aitchison similarity is computed via the centered log-ratio (CLR) transform:
\begin{equation}
  s_\text{Aitchison}(\bm{y}_A, \bm{y}_B) = \exp\!\bigl(-\|\text{clr}(\bm{y}_A) - \text{clr}(\bm{y}_B)\|_2 / 2\bigr),
\end{equation}
which respects the compositional nature of cell type fractions, unlike cosine similarity which treats the simplex as Euclidean space.

Conditions with Aitchison similarity below a threshold (default 0.5, corresponding to an Aitchison distance of approximately 1.4) are flagged for manual review, as low similarity between nominally identical conditions across screens suggests substantial batch effects. Flagged conditions may be excluded from GP training or downweighted. In the current dataset, the primary screen and SAG screen share two near-overlapping conditions (SAG\_250nM and SAG\_1uM, which closely match primary-screen conditions); these were excluded from training on the basis of this QC analysis, retaining only the two non-redundant SAG conditions (SAG\_50nM and SAG\_2uM).

\subsection*{Brain region targeting and anterior-posterior axis encoding}

The 9 brain regions annotated in \hnoca{} span the anterior-posterior (A-P) axis of the developing brain. We map each region to a continuous A-P position based on its anatomical location: Dorsal telencephalon ($a{=}0.0$), Ventral telencephalon ($a{=}0.1$), Hypothalamus ($a{=}0.2$), Thalamus ($a{=}0.3$), Dorsal midbrain ($a{=}0.5$), Ventral midbrain ($a{=}0.5$), Cerebellum ($a{=}0.7$), Pons ($a{=}0.85$), and Medulla ($a{=}1.0$). Dorsal and ventral midbrain share the same A-P position because they occupy the same rostro-caudal level (mesencephalon) and differ only along the dorso-ventral axis. These positions reflect the relative anterior-to-posterior ordering of brain structures during early fetal development.

Given a target A-P position $p$ and width parameter $\sigma$, the target cell type profile $\bm{t}$ is constructed by Gaussian-weighted combination of per-region reference profiles $\bm{r}_j$ from the HNOCA atlas:
\begin{equation}
  \bm{t}(p, \sigma) = \mathcal{C}\!\left[\sum_{j=1}^{9} \exp\!\left(-\frac{(a_j - p)^2}{2\sigma^2}\right) \bm{r}_j\right],
\end{equation}
where $\mathcal{C}$ denotes closure (normalization to the simplex) and $\bm{r}_j$ is the cell type composition vector for brain region $j$ in the reference atlas. Cell types annotated as off-target (pluripotent stem cells, mesenchymal, endothelial, microglia, neural crest) are assigned zero weight regardless of the A-P targeting. The default width $\sigma{=}0.15$ covers approximately 2--3 adjacent regions, producing a smoothly varying target that avoids hard boundaries between anatomical regions.

\subsection*{Target profile refinement}

After each experimental round, the target composition can be refined to better reflect the empirical relationship between cell type proportions and organoid quality. For each cell type $k$, we compute the Pearson correlation $\rho_k$ between its fraction across all observed conditions and the composite fidelity score. A ``learned'' target profile is then constructed as:
\begin{equation}
  t_k^\text{learned} = \frac{\exp(\rho_k / \tau)}{\sum_{j} \exp(\rho_j / \tau)},
\end{equation}
where $\tau$ is a temperature parameter (default $\tau{=}1.0$) controlling the sharpness of the softmax. Cell types with strong positive correlations to fidelity receive higher target weights, while those that are negatively correlated are suppressed.

The refined target profile is a convex combination of the original reference-derived target and the learned profile:
\begin{equation}
  \bm{t}_\text{refined} = (1 - \lambda)\,\bm{t}_\text{original} + \lambda\,\bm{t}_\text{learned},
\end{equation}
with learning rate $\lambda{=}0.3$. This conservative interpolation ensures that the optimization objective does not drift excessively from the biological prior in early rounds when the number of observed conditions is small and correlation estimates are noisy. The learning rate can be increased in later rounds as the empirical signal strengthens.

\subsection*{Adaptive bounds computation}

The GP surrogate model requires bounded input domains for acquisition function optimization. Rather than using fixed literature-derived bounds, which may include biologically implausible concentrations or irrelevant dimensions, we compute bounds adaptively from the training data. The \texttt{\_compute\_active\_bounds()} procedure first identifies and drops zero-variance dimensions---morphogens that were held constant or absent across all training conditions. In the current dataset, this removes 5 of 24 dimensions (SHH, XAV939, Dorsomorphin, purmorphamine, cyclopamine).

For the remaining active dimensions, bounds are set as:
\begin{equation}
  \text{lb}_j = \begin{cases} 0 & \text{if } j \neq \log d_\text{harvest} \\ \min(\bm{x}_j) & \text{if } j = \log d_\text{harvest} \end{cases}, \quad
  \text{ub}_j = \min\!\bigl(L_j,\; \max(\bm{x}_j) + 0.05 \cdot r_j\bigr),
\end{equation}
where $r_j = \max(\bm{x}_j) - \min(\bm{x}_j)$ is the observed range, $\bm{x}_j$ is the vector of observed values for dimension $j$, and $L_j$ is the literature maximum concentration for that morphogen. The 5\% padding on the upper bound allows the GP to recommend concentrations slightly beyond the observed range while the clamping to literature maxima prevents extrapolation into biologically implausible regimes (e.g., cytotoxic concentrations). Lower bounds are set to zero for all morphogen concentrations because negative concentrations are physically meaningless; the sole exception is the log harvest day dimension, whose lower bound is set to the training minimum to prevent recommending biologically implausible culture durations.

\subsection*{Gaussian process surrogate model}

Independent \texttt{SingleTaskGP} models (one per ILR coordinate) were fit using \botorch{} v0.11+ \cite{Balandat2020botorch} with Mat\'{e}rn 5/2 kernels. Input normalization (\texttt{Normalize}) and output standardization (\texttt{Standardize}) were applied. Hyperparameters were optimized by maximizing the marginal log-likelihood. The dimensionality-scaled lengthscale prior \cite{Hvarfner2024} places a log-normal distribution with mode $\sim\!\sqrt{d}$ on each lengthscale. SAASBO \cite{Eriksson2021saasbo} uses NUTS sampling (256 warmup, 128 samples, thinning 16) with half-Cauchy priors. All GP computations were performed on CPU (float64 precision).

\textbf{Additive + interaction kernel decomposition.} Following NAIAD \cite{Qin2025naiad}, the pipeline supports a structured kernel that decomposes the covariance function as $k(\bm{x}, \bm{x}') = k_\text{add}(\bm{x}, \bm{x}') + k_\text{int}(\bm{x}, \bm{x}')$. The additive component $k_\text{add}$ is a sum of $d$ independent 1D Mat\'{e}rn 5/2 kernels (one per morphogen dimension), each wrapped in a \texttt{ScaleKernel}, capturing independent main effects with $O(d)$ parameters. The interaction component $k_\text{int}$ is a full ARD Mat\'{e}rn 5/2 kernel over all $d$ dimensions, modeling higher-order morphogen--morphogen interactions. The interaction outputscale is initialized to 0.1, encoding a soft prior toward additivity: the model assumes most variation is explained by independent morphogen effects unless the data provide strong evidence of interactions. This kernel is selected by the adaptive complexity schedule when $N/d$ is in the intermediate regime, or can be explicitly requested via the \texttt{--kernel additive\_interaction} flag.

\textbf{Per-cell-type GP for interpretability.} As an alternative to fitting a single multi-output GP, the pipeline supports fitting separate \texttt{SingleTaskGP} models per ILR coordinate (or equivalently, per cell type) using a \texttt{ModelListGP} wrapper, inspired by GPerturb \cite{Xing2025gperturb}. Each sub-model receives its own Mat\'{e}rn 5/2 + ARD kernel with independent lengthscale hyperparameters. After fitting, a per-output lengthscale matrix $\Lambda \in \Real^{d \times K}$ is extracted, where entry $\Lambda_{j,k}$ is the ARD lengthscale of morphogen $j$ for output $k$. Shorter lengthscales indicate higher sensitivity, so this matrix directly answers which morphogens drive which cell types. This per-type GP mode is activated via the \texttt{--per-type-gp} flag and is restricted to the standard MAP fitting path (incompatible with SAASBO or multi-fidelity models).

\textbf{Warm-starting and convergence monitoring.} After each optimization round, GP hyperparameters (lengthscales, outputscale, noise variance, mean constant) are saved to disk in \texttt{.pt} format. When \texttt{--warm-start} is enabled, these saved hyperparameters initialize the next round's GP before MLL fitting continues to refine them. This warm-starting accelerates convergence across rounds by providing an informed starting point rather than default initialization. For \texttt{ModelListGP} models (SAASBO or per-type GP), per-sub-model state dictionaries are saved and restored individually, with dimension mismatch checks that gracefully skip warm-starting when the active dimension set changes between rounds.

Convergence is monitored via three diagnostics computed after each round \cite{Narayanan2025}: (1)~mean posterior standard deviation across the design space, estimated via 512 quasi-random Sobol evaluation points; (2)~peak acquisition function value in the recommended batch, which decays as the GP believes less improvement remains; and (3)~recommendation spread, the mean pairwise Euclidean distance between recommended morphogen vectors normalized to $[0,1]$. When both the acquisition value has decayed below 10\% of its round-1 level and the recommendation spread falls below a clustering threshold, the system suggests reducing the batch size, signaling that the optimizer is converging on a local region and further exploration has diminishing returns. All metrics are appended to a persistent CSV for trend visualization across rounds.

\subsection*{ILR transform for compositional responses}

Cell type fractions $\bm{y} \in \simplex^{D-1}$ were transformed to ILR coordinates using the Helmert sub-basis. Zero components were handled by multiplicative replacement \cite{MartinFernandez2015}: zeros replaced by $\delta = 0.65 / (100D)$, non-zero components scaled down proportionally to maintain closure. The Helmert sub-basis was chosen over the additive log-ratio (ALR) and centered log-ratio (CLR) alternatives for several reasons. ALR coordinates depend on the choice of reference component, introducing an asymmetry that can affect GP performance. CLR coordinates lie in a $(D{-}1)$-dimensional subspace of $\Real^D$ due to the zero-sum constraint, creating a singular covariance structure. The ILR transform maps compositions bijectively to $\Real^{D-1}$, producing unconstrained coordinates with a non-singular covariance---ideal for GP regression with standard kernels.

The inverse transform is applied to GP predictions before computing the acquisition function, ensuring that all recommended conditions correspond to valid compositions on the simplex. This is critical because GP predictions in ILR space are not constrained to map back to non-negative compositions; without explicit inverse transformation and closure, the optimizer could exploit numerically valid but biologically impossible predictions.

\subsection*{Acquisition function and recommendation}

The acquisition function was $q$-log Expected Improvement ($q$LogEI) \cite{Balandat2020botorch} with composition-space scalarization. For each candidate morphogen vector $\bm{x}$, the GP posterior provides a distribution over ILR coordinates, which is inverse-transformed to the simplex to obtain a predicted composition $\hat{\bm{y}}(\bm{x})$. This predicted composition is then scored against the target profile $\bm{t}$ via weighted inner product $\hat{\bm{y}}(\bm{x})^\top \bm{t}$, producing a scalar acquisition value. This composition-space scalarization ensures that the acquisition objective has direct biological meaning---similarity to the desired cell type distribution---rather than operating in the abstract ILR coordinate system where the relationship to biology is obscured.

For multi-objective settings where multiple brain regions are targeted simultaneously, $q$-log Noisy Expected Hypervolume Improvement ($q$LogNEHVI) was used with a data-driven reference point set at the 10th percentile of observed fidelity values. Each optimization round produces a 24-well plate map: 20 novel conditions selected by the acquisition function, 2 high-variance replicates (conditions with the highest GP posterior variance, providing maximal information gain for model refinement), and 2 QC duplicates of top-performing conditions from previous rounds (providing internal consistency controls).

\subsection*{Multi-fidelity GP}

When virtual data sources were available, a \texttt{SingleTaskMultiFidelityGP} was fit with a continuous fidelity parameter appended as an additional input dimension. The fidelity dimension encodes the data source reliability: real scRNA-seq experiments receive $f{=}1.0$, CellRank~2 temporal projections receive $f{=}0.5$, and CellFlow generative predictions receive $f{=}0.0$. The multi-fidelity kernel models the cross-fidelity correlation structure, learning to weight high-fidelity observations more heavily while still extracting useful signal from cheaper virtual data when it is informative.

An automatic gating mechanism evaluates the Spearman rank correlation $\rho$ between each virtual source and real data on overlapping or near-overlapping conditions before deciding whether to include that source. Sources with $\rho < 0.3$ are excluded entirely, as their predictions are too weakly correlated with real data to improve the surrogate. Sources with $\rho > 0.9$ are treated as near-perfect proxies and undergo pre-filtering only (removing outliers beyond 3 standard deviations). Sources with $0.3 \leq \rho \leq 0.9$ trigger full multi-fidelity modeling. Cost ratios control the acquisition function's preference for real versus virtual evaluations: real experiments cost 1.0, CellRank~2 projections cost 0.005, and CellFlow predictions cost 0.001, reflecting the approximately 200$\times$ and 1,000$\times$ cost advantage of virtual evaluations.

\textbf{Targeted Variance Reduction (TVR).} As an alternative to the continuous-fidelity GP, the pipeline supports Targeted Variance Reduction \cite{Fare2022tvr, McDonald2025mfbo}, which fits independent \texttt{SingleTaskGP} models for each fidelity level rather than modeling cross-fidelity correlation with a single kernel. At each candidate point, TVR selects the model with the lowest cost-scaled posterior variance: $\text{score}(f, \bm{x}) = \text{Var}_f(\bm{x}) / c_f$, where $\text{Var}_f$ is the posterior variance from the fidelity-$f$ GP and $c_f$ is the associated evaluation cost. Cheaper models are preferred when their raw variance is comparable to that of the expensive model, enabling the acquisition function to preferentially request low-cost evaluations in regions where they are sufficiently informative. TVR is activated via the \texttt{--use-tvr} flag and requires at least two fidelity levels in the training data. This approach trades off the ability to share information across fidelity levels (as the continuous-fidelity GP does) for greater robustness when the cross-fidelity correlation structure is poorly identified due to limited overlap between data sources.

\subsection*{Software}

The pipeline was implemented in Python 3.11+ using scanpy \cite{Wolf2018scanpy}, anndata \cite{Virshup2024anndata}, scvi-tools, \scarches{} \cite{Lotfollahi2022scarches}, \botorch{} \cite{Balandat2020botorch}, GPyTorch, PyTorch, moscot \cite{Klein2025moscot}, CellRank~2 \cite{Weiler2024cellrank2}, and Plotly. The test suite comprises 220 tests.

%% ================================================================
%% BACK MATTER
%% ================================================================

\begin{acknowledgements}
This work was self-funded by Engram. We thank the authors of the HNOCA, Braun fetal brain atlas, and Amin/Kelley morphogen screen for making their data publicly available. We acknowledge the developers of BoTorch, scArches, CellRank, moscot, and CellFlow for their open-source software contributions. The GP-BO pipeline source code will be made available on GitHub upon publication.
\end{acknowledgements}

\begin{contributions}
M.Y. conceived the project, designed and implemented the computational framework, and wrote the manuscript. R.Z. contributed to atlas mapping and fidelity scoring analysis. J.K. contributed to multi-fidelity integration and experimental validation design. W.S. contributed to GP model design and acquisition function development. All authors reviewed and edited the manuscript.
\end{contributions}

\begin{interests}
Engram has filed patent applications related to the computational methods described in this work.
\end{interests}

\bibliography{references}

%% ================================================================
%% FIGURES (full-width, at end)
%% ================================================================

\begin{figure*}[tbp]
\centering
\fbox{\parbox{0.90\textwidth}{\centering\vspace{2.5cm}{\large\sffamily\bfseries Figure 1: Pipeline Overview}\\\vspace{0.3cm}\textit{(Figure placeholder --- to be generated)}\vspace{2.5cm}}}
\caption{\textbf{Multi-fidelity GP-BO framework for brain organoid morphogen optimization.}
\textbf{a},~Morphogen conditions are encoded as 24-dimensional vectors spanning WNT, BMP, SHH, RA, FGF, Notch, EGF, and Activin signaling pathways plus harvest timing.
\textbf{b},~scRNA-seq data from organoid cultures are mapped onto the HNOCA reference atlas via scArches/scPoli transfer learning, producing cell type fraction vectors per condition.
\textbf{c},~Cell type fractions are ILR-transformed and used to fit a multi-output GP surrogate model.
\textbf{d},~An acquisition function in composition space selects the next 24 experiments.
\textbf{e},~Multi-fidelity data integration combines real experiments ($f{=}1.0$) with CellRank~2 temporal projections ($f{=}0.5$) and CellFlow generative predictions ($f{=}0.0$).}
\label{fig:overview}
\end{figure*}

\begin{figure*}[tbp]
\centering
\fbox{\parbox{0.90\textwidth}{\centering\vspace{2.5cm}{\large\sffamily\bfseries Figure 2: Fidelity Scoring}\\\vspace{0.3cm}\textit{(Figure placeholder --- to be generated)}\vspace{2.5cm}}}
\caption{\textbf{Two-tier fidelity scoring against the fetal brain.}
\textbf{a},~Schematic of the four sub-metrics.
\textbf{b},~Composite fidelity scores across 46 primary-screen conditions, colored by dominant brain region.
\textbf{c},~Cell type composition stacked bar plot for selected conditions.
\textbf{d},~Brain region assignment across all conditions, showing 8 distinct dominant regions.}
\label{fig:fidelity}
\end{figure*}

\begin{figure*}[tbp]
\centering
\fbox{\parbox{0.90\textwidth}{\centering\vspace{2.5cm}{\large\sffamily\bfseries Figure 3: Multi-fidelity Integration}\\\vspace{0.3cm}\textit{(Figure placeholder --- to be generated)}\vspace{2.5cm}}}
\caption{\textbf{Multi-fidelity data integration augments the GP training set.}
\textbf{a},~Data sources at three fidelity levels with associated costs.
\textbf{b},~Cross-fidelity correlation gating.
\textbf{c},~GP posterior uncertainty with and without multi-fidelity data.}
\label{fig:multifidelity}
\end{figure*}

\begin{figure*}[tbp]
\centering
\fbox{\parbox{0.90\textwidth}{\centering\vspace{2.5cm}{\large\sffamily\bfseries Figure 4: Round 1 Results}\\\vspace{0.3cm}\textit{(Figure placeholder --- to be generated)}\vspace{2.5cm}}}
\caption{\textbf{Round~1 GP-BO analysis and recommendations.}
\textbf{a},~PCA projection of the 24-dimensional morphogen space.
\textbf{b},~GP lengthscale estimates (MAP) reveal relative morphogen importance.
\textbf{c},~Acquisition function landscape showing recommended conditions.
\textbf{d},~Predicted cell type compositions for the top 6 recommended conditions.
\textbf{e},~Leaderboard of round~1 conditions ranked by composite fidelity.}
\label{fig:results}
\end{figure*}

\begin{table*}[tbp]
\centering
\caption{\textbf{The 24-dimensional morphogen search space.} All concentrations are in \si{\micro\Molar}. Protein molecular weights used for unit conversion are shown where applicable. Dims 20--23 are base media components held constant across the primary screen.}\label{tab:morphogens}
\small
\begin{tabular}{@{}clllc@{}}
\toprule
\textbf{Dim} & \textbf{Column} & \textbf{Pathway} & \textbf{Type} & \textbf{MW (kDa)} \\
\midrule
0  & CHIR99021  & WNT agonist      & Small molecule & --- \\
1  & BMP4       & BMP              & Protein        & 13.0 \\
2  & BMP7       & BMP              & Protein        & 15.7 \\
3  & SHH        & Sonic hedgehog   & Protein        & 19.6 \\
4  & SAG        & SHH agonist      & Small molecule & --- \\
5  & RA         & Retinoic acid    & Small molecule & --- \\
6  & FGF8       & FGF              & Protein        & 22.5 \\
7  & FGF2       & FGF              & Protein        & 17.2 \\
8  & FGF4       & FGF              & Protein        & 19.2 \\
9  & IWP2       & WNT inhibitor    & Small molecule & --- \\
10 & XAV939     & WNT inhibitor    & Small molecule & --- \\
11 & SB431542   & TGF-$\beta$ inh. & Small molecule & --- \\
12 & LDN193189  & BMP inhibitor    & Small molecule & --- \\
13 & DAPT       & Notch inhibitor  & Small molecule & --- \\
14 & EGF        & EGF              & Protein        & 6.2 \\
15 & Activin A  & TGF-$\beta$      & Protein        & 26.0 \\
16 & Dorsomorphin & BMP inhibitor  & Small molecule & --- \\
17 & Purmorphamine & SHH agonist   & Small molecule & --- \\
18 & Cyclopamine & SHH antagonist  & Small molecule & --- \\
19 & $\log(d_\text{harvest})$ & Time & --- & --- \\
20 & BDNF       & Neurotrophin     & Protein        & 13.5 \\
21 & NT-3       & Neurotrophin     & Protein        & 13.6 \\
22 & cAMP       & cAMP signaling   & Small molecule & --- \\
23 & Ascorbic acid & Antioxidant   & Small molecule & --- \\
\bottomrule
\end{tabular}
\end{table*}

%% ================================================================
%% SUPPLEMENTARY INFORMATION
%% ================================================================
\clearpage
\onecolumn
\setcounter{figure}{0}
\setcounter{table}{0}
\renewcommand{\thefigure}{S\arabic{figure}}
\renewcommand{\thetable}{S\arabic{table}}

\section*{Supplementary Information}

\subsection*{Supplementary Table S1: Fidelity scores for all 46 primary-screen conditions}

\begin{table}[H]
\centering
\caption{\textbf{Composite fidelity scores and sub-metrics for representative conditions.} Top 10 and bottom 5 conditions from the 46 primary-screen conditions are shown. All conditions show 14 detected cell types at HNOCA level-2. Off-target fractions are uniformly low ($<5\%$), indicating successful neural induction across all protocols.}\label{tab:fidelity-supp}
\small
\begin{tabular}{@{}lrccccc@{}}
\toprule
\textbf{Condition} & \textbf{$n$ cells} & \textbf{Region} & \textbf{On-target} & \textbf{Off-target} & \textbf{RSS} & \textbf{Fidelity} \\
\midrule
\multicolumn{7}{l}{\textit{Top 10 conditions}} \\
IWP2 & 858 & Dorsal tel. & 0.969 & 0.000 & 1.000 & 0.979 \\
IWP2 switch CHIR & 557 & Dorsal tel. & 0.932 & 0.000 & 0.999 & 0.978 \\
FGF4 & 1,235 & Dorsal tel. & 0.938 & 0.000 & 0.979 & 0.977 \\
FGF8 & 944 & Dorsal tel. & 0.974 & 0.000 & 0.971 & 0.972 \\
CHIR-d16-21 & 1,670 & Dorsal tel. & 0.938 & 0.000 & 0.999 & 0.968 \\
DAPT & 903 & Dorsal tel. & 0.951 & 0.000 & 0.932 & 0.963 \\
LDN & 1,236 & Dorsal tel. & 0.970 & 0.000 & 0.996 & 0.942 \\
I/Activin/DAPT/SR11 & 1,423 & Dorsal tel. & 0.784 & 0.001 & 0.976 & 0.937 \\
CHIR1.5 & 1,154 & Thalamus & 0.654 & 0.000 & 0.917 & 0.848 \\
C/S/R/E/FGF2/D & 492 & Dorsal mid. & 0.524 & 0.002 & 0.844 & 0.826 \\
\midrule
\multicolumn{7}{l}{\textit{Bottom 5 conditions}} \\
RA100 & 190 & Unspecific & 0.384 & 0.016 & 0.324 & 0.560 \\
SAG-CHIRd10-21 & 368 & Unspecific & 0.432 & 0.030 & 0.214 & 0.557 \\
CHIR3-SAG1000 & 499 & Unspecific & 0.443 & 0.048 & 0.413 & 0.555 \\
C/S/D/FGF4 & 126 & Dorsal mid. & 0.246 & 0.000 & 0.483 & 0.543 \\
\bottomrule
\end{tabular}
\end{table}

\subsection*{Supplementary Table S2: GP-BO recommended conditions for Round 2}

\begin{table}[H]
\centering
\caption{\textbf{Round~1 GP-BO recommendations.} The acquisition function recommended 4 novel conditions and 2 replicates. Concentrations in \si{\micro\Molar}; harvest day shown as $\log(d)$.}\label{tab:recommendations-supp}
\small
\begin{tabular}{@{}lccccl@{}}
\toprule
\textbf{Well} & \textbf{CHIR} & \textbf{BMP4} & \textbf{SAG} & \textbf{$\log(d)$} & \textbf{Type} \\
\midrule
A1 & 0.212 & 0.0008 & 0.086 & 3.15 & Novel \\
A2 & 2.172 & 0.0032 & 0.591 & 2.30 & Novel \\
A3 & 2.489 & 0.0003 & 0.685 & 3.15 & Novel \\
A4 & 0.827 & 0.0013 & 0.382 & 4.08 & Novel \\
A5 & 2.422 & 0.0036 & 0.318 & 0.50 & Replicate \\
A6 & 2.789 & 0.0032 & 0.634 & 3.92 & Replicate \\
\bottomrule
\end{tabular}
\end{table}

\subsection*{Supplementary Figure S1--S4 (Placeholders)}

\begin{figure}[H]
\centering
\fbox{\parbox{0.85\textwidth}{\centering\vspace{2cm}{\sffamily\bfseries Supplementary Figure S1}\\\textit{Cell type composition stacked bar plots for all 46 conditions.}\vspace{2cm}}}
\caption{\textbf{Cell type composition across all primary-screen conditions.} Stacked bar plots showing the fraction of each of 14 HNOCA level-2 cell types per condition, ordered by composite fidelity score.}
\label{fig:supp-composition}
\end{figure}

\begin{figure}[H]
\centering
\fbox{\parbox{0.85\textwidth}{\centering\vspace{2cm}{\sffamily\bfseries Supplementary Figure S2}\\\textit{GP posterior mean and uncertainty for each morphogen dimension.}\vspace{2cm}}}
\caption{\textbf{GP surrogate model diagnostics.} Posterior mean and 95\% credible interval for each active morphogen dimension, showing how the GP interpolates between observed conditions.}
\label{fig:supp-gp}
\end{figure}

\begin{figure}[H]
\centering
\fbox{\parbox{0.85\textwidth}{\centering\vspace{2cm}{\sffamily\bfseries Supplementary Figure S3}\\\textit{Brain region assignment UMAP.}\vspace{2cm}}}
\caption{\textbf{UMAP visualization of mapped organoid cells colored by brain region assignment.} Cells from all 46 primary-screen conditions projected into the HNOCA latent space, colored by transferred brain region label.}
\label{fig:supp-umap}
\end{figure}

\begin{figure}[H]
\centering
\fbox{\parbox{0.85\textwidth}{\centering\vspace{2cm}{\sffamily\bfseries Supplementary Figure S4}\\\textit{Algorithm pseudocode for the GP-BO optimization loop.}\vspace{2cm}}}
\caption{\textbf{GP-BO optimization loop pseudocode.} Complete algorithm showing data loading, ILR transform, GP fitting, acquisition optimization, and plate map generation for each round.}
\label{fig:supp-algo}
\end{figure}

\end{document}
